% ==============================================================================
% Appendix D: Hardware Silicon Dissection & Architecture Comparison:
%             NVIDIA AD104 vs. Huawei Ascend 910B
% File: academic/paper/arxiv/sections/sec_appendix_d_silicon.tex
% ==============================================================================

\section{Hardware Silicon Dissection and Architecture Comparison: NVIDIA AD104 vs. Huawei Ascend 910B}
\label{sec:appendix_silicon_dissection}

The theoretical computational complexity of sequence modeling algorithms---conventionally expressed through asymptotic bounds such as $\mathcal{O}(N)$ for linear recurrences versus $\mathcal{O}(N^2)$ for canonical softmax self-attention---fails to predict actual wall-clock execution latency, memory subsystem saturation, and energy expenditure on physical silicon. The realization of algorithmic speedups in large language models (LLMs) depends intrinsically on the strict mathematical alignment of tensor traversal patterns with the physical memory hierarchy of the target processor: thread-private register files, on-chip scratchpad SRAM and shared memory, multi-megabyte crossbar L2 caches, and high-bandwidth off-chip global memory (GDDR6X and HBM2e).

Within the \neuralops{} framework, understanding the microarchitectural root causes of performance cliffs, register spilling, shared memory bank collisions, and DMA double-buffering stalls is necessary to guide automated kernel selection and in-place attention surgery. To establish rigorous empirical and architectural grounding, this Appendix provides an exhaustive microarchitectural dissection and comparative analysis of the two primary heterogeneous compute substrates evaluated in our research:
\begin{enumerate}
    \item \textbf{NVIDIA Ada Lovelace (AD104 Silicon, GeForce RTX 4070 / RTX 4070 Ti, SM 8.9)}: Representing the state-of-the-art desktop, edge, and workstation computing paradigm. It features an integrated SIMT execution pipeline, 4th-generation Tensor Cores with hardware-accelerated FP8 Transformer Engine support, structured 2:4 sparsity, an asynchronous copy engine (\texttt{cp.async}), and an enlarged $48\text{ MB}$ L2 cache crossbar designed to mitigate global memory bandwidth bottlenecks.
    \item \textbf{Huawei Ascend 910B (DaVinci v2 Architecture, CANN Platform)}: Representing the disaggregated, domain-specific accelerator paradigm. It physically decouples matrix operations from elementwise execution via independent Cube Units ($16 \times 16 \times 16$ 3D fractal matrix MAC engines) and Vector Units (wide SIMD vector processors), supported by an explicit, software-managed $256\text{ KB}$ Unified Buffer (UB) scratchpad SRAM, multi-stage Memory Transfer Engines (MTE2/MTE3), and strict 32-byte/512-byte memory alignment contracts.
\end{enumerate}

Additionally, to provide broader context across modern datacenter architectures, comparative microarchitectural parameters for the NVIDIA Ampere (A100 SXM4, SM 8.0) and Hopper (H100 SXM5, SM 9.0a) architectures are incorporated throughout this analysis.

% ==============================================================================
% D.1 NVIDIA Ada Lovelace (AD104) Microarchitecture Dissection
% ==============================================================================
\subsection{NVIDIA Ada Lovelace (AD104) Microarchitecture Dissection}
\label{sec:appendix_nvidia_ada_dissection}

\subsubsection{Physical Silicon Floorplan and Compute Topology}
\label{sec:appendix_ada_floorplan}

The NVIDIA GeForce RTX 4070 and RTX 4070 Ti GPUs are fabricated on the custom TSMC 4N process node (an optimized $5\text{ nm}$ lithographic node). The monolithic AD104 die integrates \textbf{$35.8$ billion transistors} within a physical die area of \textbf{$294.5\text{ mm}^2$}, achieving an active transistor density of approximately $121.6\text{ million transistors/mm}^2$.

The silicon floorplan of the full AD104 die is organized into 5 Graphics Processing Clusters (GPCs). Each fully enabled GPC houses 6 Texture Processing Clusters (TPCs), with each TPC containing 2 Streaming Multiprocessors (SMs), yielding a full physical capacity of \textbf{60 Streaming Multiprocessors (SMs)}. In the primary experimental testbed utilized throughout this monograph:
\begin{itemize}
    \item \textbf{GeForce RTX 4070 (AD104, SM 8.9)}: Configured with 5,888 FP32 CUDA cores across 46 SMs, 184 4th-generation Tensor Cores, operating at a base clock of $1,695\text{ MHz}$ (dense base compute of $141.6\text{ TFLOPS}$ FP16) and a boost clock of $2,175\text{ MHz}$ (peak dense FP16 tensor throughput of $160.2\text{ TFLOPS}$, and $320.4\text{ TFLOPS}$ with 2:4 structured sparsity; non-tensor vector FP32 compute ceiling of $70.8\text{ TFLOPS}$), backed by 12~GB of GDDR6X memory @ $504.0\text{ GB/s}$ ($I^* = 317.9\text{ FLOP/Byte}$).
    \item \textbf{GeForce RTX 4070 Ti (AD104-400)}: Configured with all 60 physical SMs fully enabled, exposing 7,680 FP32 CUDA cores, 240 4th-generation Tensor Cores, and 60 RT cores, operating at a boost clock of $2,610\text{ MHz}$.
\end{itemize}

\subsubsection{Internal SM Partitioning and Sub-Core Microarchitecture}
\label{sec:appendix_ada_sm_partitioning}

Each Ada Lovelace Streaming Multiprocessor (Compute Capability \texttt{sm\_8.9}) is architecturally partitioned into four functionally independent processing sub-cores (often designated as SMSlices or Processing Blocks), as illustrated in Figure~\ref{fig:ad104_sm_pipeline}. Each sub-core operates on a distinct quarter of the thread warps assigned to the SM and contains:
\begin{enumerate}
    \item \textbf{Warp Scheduler and Instruction Dispatch Unit}: One warp scheduler capable of issuing up to two instructions per clock cycle to its execution units, supported by an instruction cache and branch prediction logic.
    \item \textbf{Sub-Core Register File}: A 16,384-entry ($16\text{K}$) 32-bit physical register file slice, providing $64\text{ KB}$ of ultra-fast SRAM storage per sub-core. Combined across the four sub-cores, each SM hosts \textbf{65,536 (64K) 32-bit registers}, totaling \textbf{$256\text{ KB}$ of register SRAM per SM}. Across the 60 SMs of the full AD104 die, the aggregate register file bandwidth exceeds $54.2\text{ TB/s}$, with a cumulative on-chip register capacity of \textbf{$15.36\text{ MB}$} ($11.78\text{ MB}$ across 46 SMs).
    \item \textbf{Dual-Issue Execution Datapaths}: Each sub-core contains 32 FP32 ALUs (128 FP32 ALUs per SM) and 16 INT32 ALUs (64 INT32 ALUs per SM). The Ada Lovelace sub-core features an independent concurrent execution pipeline that allows the warp scheduler to simultaneously issue one FP32 instruction and one INT32 instruction in the same cycle. In the \neuralops{} Level 1 CUDA kernels (e.g., \texttt{selective\_scan} and \texttt{chunked\_delta\_gate\_scan}), this dual-issue capability is exploited to overlap 64-bit tensor stride indexing arithmetic with floating-point multiply-accumulate computations without incurring stall cycles.
    \item \textbf{Special Function Units (SFUs)}: 4 SFUs per sub-core (16 SFUs per SM) dedicated to computing hardware-approximated transcendental and non-linear functions, including $\sin(x)$, $\cos(x)$, $\exp_2(x)$, and $\log_2(x)$. These units provide the physical throughput required for the trigonometric coordinate mappings in \texttt{cosformer\_linear\_attention} and rotary position embeddings (\texttt{fused\_rotary\_emb}).
    \item \textbf{4th-Generation Tensor Core}: 1 Tensor Core per sub-core (4 Tensor Cores per SM), delivering hardware-accelerated matrix-multiply-accumulate (MMA) operations for FP16, BF16, TF32, and native FP8 numerical formats.
\end{enumerate}

% ------------------------------------------------------------------------------
% TikZ Figure D.1: AD104 SM Pipeline Schematic
% ------------------------------------------------------------------------------
\begin{figure*}[t]
\centering
\resizebox{0.98\linewidth}{!}{%
\begin{tikzpicture}[
    scale=0.88, every node/.style={transform shape},
    box/.style={draw=black!70, rounded corners=2pt, line width=0.8pt, align=center, font=\footnotesize},
    smbox/.style={draw=blue!80!black, fill=blue!5, rounded corners=4pt, line width=1.2pt},
    subcore/.style={draw=blue!70!black, fill=white, rounded corners=3pt, line width=0.9pt},
    membox/.style={draw=green!60!black, fill=green!8, rounded corners=2pt, line width=0.8pt, font=\footnotesize},
    cachebox/.style={draw=orange!80!black, fill=orange!8, rounded corners=3pt, line width=1pt, font=\footnotesize},
    drambox/.style={draw=red!70!black, fill=red!8, rounded corners=3pt, line width=1.1pt, font=\small\bfseries},
    arr/.style={-{Stealth[length=5pt,width=4pt]}, line width=1pt, draw=black!75},
    darr/.style={{Stealth[length=5pt,width=4pt]}-{Stealth[length=5pt,width=4pt]}, line width=1.1pt, draw=black!75}
]

% Main SM Box
\draw[smbox] (-0.5, 3.2) rectangle (15.5, 10.8);
\node[anchor=north west, font=\normalsize\bfseries, text=blue!90!black] at (-0.3, 10.65) {NVIDIA Ada Lovelace Streaming Multiprocessor (SM 8.9) --- 60 SMs in AD104};

% 4 Processing Sub-Cores (SMSlices)
\foreach \i in {0, 1, 2, 3} {
    \pgfmathsetmacro{\xoffset}{\i * 3.8 + 0.0}
    \draw[subcore] (\xoffset, 5.2) rectangle (\xoffset + 3.5, 10.1);
    \node[anchor=north, font=\scriptsize\bfseries, text=blue!80!black] at (\xoffset + 1.75, 10.0) {SMSlice \i{} (Sub-Core)};
    
    % Warp Scheduler & Dispatch
    \node[box, fill=purple!10, text width=3.1cm] at (\xoffset + 1.75, 9.25) {Warp Scheduler \& Dispatch\\{\tiny (1 Warp/Cycle Dispatch, 32 Threads)}};
    
    % Register File Slice
    \node[box, fill=yellow!15, text width=3.1cm] at (\xoffset + 1.75, 8.2) {Register File Slice\\{\tiny 16K $\times$ 32-bit Regs (64 KB SRAM)}};
    
    % Execution Units
    \node[box, fill=blue!10, text width=3.1cm] at (\xoffset + 1.75, 7.15) {32 $\times$ FP32 ALUs / 16 $\times$ INT32\\{\tiny Concurrent Dual-Issue Math Execution}};
    
    % Tensor Core & SFU
    \node[box, fill=teal!15, text width=3.1cm] at (\xoffset + 1.75, 5.95) {4th-Gen Tensor Core + 4 $\times$ SFU\\{\tiny FP8/FP16/BF16/TF32 MMA + 2:4 Sparse}};
}

% Unified L1 / Shared Memory
\draw[membox, fill=green!12] (0.0, 3.5) rectangle (15.0, 4.7);
\node[font=\footnotesize\bfseries, text=green!50!black] at (7.5, 4.35) {Configurable Unified L1 Data Cache / Shared Memory (SRAM)};
\node[font=\scriptsize, text=black!80] at (7.5, 3.9) {Total Capacity: 128 KB per SM \quad$\vert$\quad Max User Shared Memory: 99 KB per Threadblock \quad$\vert$\quad 32 Independent 4-Byte Banks};

% Interconnect SM to L1/Smem
\draw[darr] (7.5, 5.2) -- (7.5, 4.7);

% L2 Cache
\draw[cachebox] (1.0, 1.6) rectangle (14.0, 2.7);
\node[font=\small\bfseries, text=orange!70!black] at (7.5, 2.35) {Massive Crossbar L2 Cache Subsystem};
\node[font=\scriptsize, text=black!80] at (7.5, 1.95) {Capacity: 48 MB (RTX 4070 Ti) / 36 MB (RTX 4070) \quad$\vert$\quad Sustained Bandwidth: $\ge 1.8\text{ TB/s}$ (Crossbar Interconnect)};

% Interconnect Smem to L2
\draw[darr] (7.5, 3.5) -- node[right, font=\tiny\bfseries, text=black!70] {\texttt{cp.async} DMA Bypass} (7.5, 2.7);

% DRAM / GDDR6X
\draw[drambox] (2.5, -0.1) rectangle (12.5, 0.9);
\node[font=\normalsize\bfseries, text=red!80!black] at (7.5, 0.55) {Global Memory (GDDR6X DRAM Subsystem)};
\node[font=\scriptsize, text=black!80] at (7.5, 0.2) {12 GB GDDR6X on 192-bit Bus \quad$\vert$\quad Theoretical Bandwidth: 504.0 GB/s \quad$\vert$\quad Sustained STREAM: 255.8 GB/s};

% Interconnect L2 to DRAM
\draw[darr] (7.5, 1.6) -- (7.5, 0.9);

\end{tikzpicture}%
}
\caption[NVIDIA Ada Lovelace (AD104, SM 8.9) Microarchitecture and Memory Subsystem]{\textbf{NVIDIA Ada Lovelace (AD104, SM 8.9) Microarchitecture and Memory Subsystem Hierarchy.} The processor partitions each SM into four independent processing sub-cores (SMSlices), each containing a 16K 32-bit register file slice, concurrent dual-issue FP32/INT32 execution units, a 4th-generation Tensor Core, and Special Function Units. A unified $128\text{ KB}$ L1/Shared Memory pool interfaces via high-speed crossbar with the enlarged $48\text{ MB}$ L2 cache, isolating execution pipelines from the $504.0\text{ GB/s}$ GDDR6X DRAM bus.}
\label{fig:ad104_sm_pipeline}
\end{figure*}

\subsubsection{Register File Allocation Dynamics and the Occupancy Cliff}
\label{sec:appendix_ada_register_allocation}

Because the physical register file is shared across all concurrent warps allocated to an SM, the per-thread register consumption dictated by the compiler (\texttt{ptxas}) acts as the primary hardware gate governing thread-level parallelism and active warp residency. The hardware execution engine organizes threads into warps of $32$ threads. The theoretical maximum active warp capacity of an Ada Lovelace SM is \textbf{48 active warps} (1,536 threads per SM).

The upper bound on concurrently schedulable warps is strictly constrained by the total register allocation:
\begin{equation}
    \text{Active Warps per SM} = \min \left( \left\lfloor \frac{65,536}{\text{Regs per Thread} \times 32} \right\rfloor, \, 48 \right).
    \label{eq:active_warps_formula}
\end{equation}

Table~\ref{tab:register_occupancy_breakdown} details the non-linear relationship between compiler register allocation and physical hardware occupancy on the AD104 architecture.

\begin{table}[t]
\centering
\small
\caption[Register Allocation vs Theoretical SM Occupancy]{\textbf{Register File Allocation vs. Theoretical SM Occupancy on NVIDIA Ada Lovelace (SM 8.9).} Hardware limits: 65,536 physical 32-bit registers per SM; maximum 48 active warps (1,536 threads); maximum 255 registers per individual thread.}
\label{tab:register_occupancy_breakdown}
\setlength{\tabcolsep}{6pt}
\renewcommand{\arraystretch}{1.12}
\resizebox{\columnwidth}{!}{%
\begin{tabular}{cccccc}
\toprule
\textbf{Regs / Thread} & \textbf{Regs / Warp} & \textbf{Max Active Warps} & \textbf{Active Threads} & \textbf{SM Occupancy} & \textbf{Latency Tolerance Window} \\
\midrule
$\le 32$ & $1,024$ & 48 & 1,536 & \textbf{100.0\%} & Maximum latency hiding (Full pipelining) \\
$\le 40$ & $1,280$ & 40 & 1,280 & \textbf{83.3\%}  & High tolerance (Minor warp starvation) \\
$\le 48$ & $1,536$ & 32 & 1,024 & \textbf{66.7\%}  & Standard tolerance (Moderate bandwidth bounds) \\
$\le 64$ & $2,048$ & 32 & 1,024 & \textbf{66.7\%}  & Standard compute baseline \\
$\le 80$ & $2,560$ & 24 & 768   & \textbf{50.0\%}  & \textbf{Critical Occupancy Cliff} \\
$\le 104$ & $3,328$ & 19 & 608   & \textbf{39.6\%}  & Severe pipeline stall vulnerability \\
$\le 128$ & $4,096$ & 16 & 512   & \textbf{33.3\%}  & Memory stall bound; requires unrolling \\
$\ge 256$ & ---     & --- & ---  & ---              & \textbf{Fatal Local Memory Spilling} \\
\bottomrule
\end{tabular}%
}
\end{table}

The transition from $64$ to $80$ registers per thread induces the \textbf{Occupancy Cliff}: theoretical SM occupancy abruptly drops from $66.7\%$ to $50.0\%$. In memory-bandwidth-bound linear attention kernels, this reduction in active warps prevents the warp scheduler from finding eligible warps to issue during long-latency DRAM access phases (which require 200--400 clock cycles), leaving execution ALUs starved of instructions.

If an unrolled kernel exceeds the absolute hardware limit of \textbf{255 registers per thread}, the \texttt{ptxas} compiler forcibly spills the excess registers into \textbf{Local Memory}. Spilled register accesses are converted into \texttt{st.local} and \texttt{ld.local} PTX instructions. Although backed by the unified L1/L2 cache hierarchy, local memory spills incur significant address translation overheads, trigger L1 eviction thrashing, and introduce full memory-pipeline latency stalls ($200\text{--}400$ cycles upon cache misses), degrading computational throughput by up to $15\times$ to $20\times$. 

In just-in-time (JIT) compilation under OpenAI Triton, register spilling presents an acute hazard: aggressive loop unrolling across wide 2D tile accumulators ($M_{\text{tile}} \times N_{\text{tile}} \ge 128 \times 128$) and multi-stage software pipelines ($N_{\text{stages}} \ge 4$) cause LLVM's live-range analyzer to expand register allocations past the 255-register hardware ceiling. To prevent catastrophic spilling, \neuralops{} enforces launch bound directives across all Level 1 CUDA kernels:
\begin{lstlisting}[language=C++]
__global__ void __launch_bounds__(128, 2)
selective_scan_fwd_kernel(...) {
    // Restricts compiler register allocation to <= 64 registers/thread
}
\end{lstlisting}
The attribute \texttt{\_\_launch\_bounds\_\_(128, 2)} instructs \texttt{ptxas} to guarantee that at least two 128-thread blocks can execute simultaneously per SM, strictly capping register allocation at $\lfloor 65536 / (2 \times 128) \rfloor = 256$ registers per block ($64$ registers per thread), thereby preventing local memory spills. Complementarily, the Triton autotuning engine dynamically evaluates register footprints across candidate configurations, automatically pruning configurations that exceed 64 registers per thread.

\subsubsection{Shared Memory Banking and Prime Stride Conflict Elimination}
\label{sec:appendix_ada_shared_memory_banking}

The on-chip SRAM in each Ada Lovelace SM is configured as a unified pool of \textbf{$128\text{ KB}$}, dynamically partitioned between hardware L1 data cache and software-managed shared memory (\texttt{smem}). Under driver configuration, the maximum user-accessible shared memory per Cooperative Thread Array (CTA / threadblock) is capped at \textbf{$99\text{ KB}$}, with the remaining $29\text{ KB}$ permanently reserved for L1 cache tag structures and hardware data staging.

The physical shared memory array is organized into \textbf{32 independent memory banks}, each with a word width of \textbf{4 bytes (32 bits)}. Bank addresses are assigned cyclically:
\begin{equation}
    \text{Bank Index} = \left( \frac{\text{Byte Address}}{4} \right) \bmod 32.
\end{equation}
In a given clock cycle, each bank can service exactly one 32-bit read or write request. When multiple threads within the same warp simultaneously request distinct memory addresses that map to the identical bank index, a \textbf{Shared Memory Bank Conflict} occurs. The hardware arbiter must serialize the conflicting accesses:
\begin{equation}
    \text{Effective Latency} = \text{Base Latency} \times N_{\text{way}},
\end{equation}
where $N_{\text{way}} \in \{1, 2, 4, 8, 16, 32\}$ denotes the conflict degree. A 32-way bank conflict serializes the access across 32 consecutive cycles, reducing effective shared memory throughput by $96.9\%$.

In autoregressive linear attention and low-rank state-space kernels, recurrent hidden state matrices $S_t \in \R^{d \times d}$ are commonly stored in shared memory with power-of-two dimensions (e.g., $d = 64$ or $d = 128$). When a warp reads column slices of $S_t$ across threads ($S_{t}[k, i]$ for thread $i$), the access stride is an exact multiple of 32:
\begin{equation}
    \text{Address}(k, i) = \text{Base} + (k \cdot d + i) \times 4 \implies \text{Bank}(k, i) = (k \cdot d + i) \bmod 32.
\end{equation}
If $d \bmod 32 = 0$, every thread in column-traversal operations maps to the identical bank, causing a maximal 32-way conflict.

To resolve this bottleneck, \neuralops{} kernels implement \textbf{Prime Stride Padding}. For a 2D tile accumulator of shape $[M \times K]$ allocated in shared memory, padding the inner dimension by $+1$ word ($\text{Stride}_{\text{Shared}} = K + 1$) guarantees that the stride is coprime with the 32 physical memory banks:
\begin{equation}
    \text{Stride}_{\text{Shared}} = K + 1, \quad \text{where } \gcd(K + 1, \; 32) = 1.
\end{equation}
For a standard $32 \times 32$ state tile of 32-bit floats ($K=32$), setting $\text{Stride} = 33$ results in the physical bank mapping for thread $r \in [0, 31]$ accessing column $c$:
\begin{equation}
    \text{Bank ID}(r, c) = (r \times 33 + c) \bmod 32 = (r + c) \bmod 32.
\end{equation}
Because $\gcd(33, 32) = 1$, the linear congruence $r \mapsto (r + c) \bmod 32$ forms a complete bijective permutation on $\mathbb{Z}_{32}$. When all 32 threads in a warp traverse column $c$ across rows $r \in [0, 31]$, each lane maps to a distinct physical bank without collisions, achieving $100\%$ conflict-free single-cycle SRAM access ($128\text{ Bytes/cycle per SM}$). The microarchitectural data layout is encapsulated in \Cref{lst:appendix_cuda_tile}:
\begin{lstlisting}[language=C++, caption={\textbf{Bank-Conflict-Free Shared Memory Tile Layout with Prime Stride Padding (\texttt{stride = TILE\_K + 1})}}, label={lst:appendix_cuda_tile}]
template <int TILE_M = 32, int TILE_K = 32>
struct alignas(16) ConflictFreeSharedTile {
    // Prime stride padding: gcd(TILE_K + 1, 32) = gcd(33, 32) = 1 enforces a bijection
    static constexpr int STRIDE_K = TILE_K + 1;
    float tile[TILE_M][STRIDE_K];

    __device__ __forceinline__ float& at(int r, int c) {
        return tile[r][c]; // Bank ID = (r * STRIDE_K + c) % 32 = (r + c) % 32
    }
};
\end{lstlisting}

\subsubsection{The L2 Cache Revolution and Spatial Recurrence Locality}
\label{sec:appendix_ada_l2_revolution}

The most consequential microarchitectural departure of the Ada Lovelace generation is the massive expansion of the on-chip Level 2 (L2) cache. While the preceding Ampere GA104 die (GeForce RTX 3070) integrated a modest $4\text{ MB}$ L2 cache, the AD104 die integrates \textbf{$48\text{ MB}$ of crossbar L2 cache} on the RTX 4070 Ti ($36\text{ MB}$ on the RTX 4070)---a \textbf{$9\times$ to $12\times$ capacity increase}.

The AD104 L2 cache is interfaced via a high-bandwidth internal crossbar network delivering over \textbf{$1.8\text{ TB/s}$ of sustained internal bandwidth}, compared to the $504.0\text{ GB/s}$ external GDDR6X bus. This architectural shift alters the operational roofline for recurrent sequence modeling:
\begin{itemize}
    \item \textbf{Recurrent State Retention}: In chunked linear attention algorithms (\texttt{chunk\_kda}, \texttt{gated\_deltanet}), the sequence recurrence state tensor for batch size $B=2$, heads $H=12$, and state dimension $d=64$ requires:
    \begin{equation}
        \text{Footprint}(S_t) = B \times H \times d \times d \times \text{sizeof}(\text{FP32}) = 2 \times 12 \times 64 \times 64 \times 4\text{ bytes} = 393.2\text{ KB}.
    \end{equation}
    On Ampere architectures, concurrent kernel execution and inter-chunk state transfers regularly evict $S_t$ to off-chip DRAM. On Ada Lovelace, the $36\text{--}48\text{ MB}$ L2 capacity enables the entire recurrent state of all heads and batches to remain resident in on-chip SRAM/L2 across chunk boundaries.
    \item \textbf{Block Swizzling Optimization}: By applying Morton space-filling curves (Z-order swizzling) to the threadblock grid scheduler:
    \begin{equation}
        \text{GridID}(b, h) = \text{InterleaveBits}(b, h),
    \end{equation}
    \neuralops{} kernels maximize the temporal and spatial cache locality of key-value projection matrices, achieving an L2 cache hit rate exceeding $88.4\%$ during sequence execution.
\end{itemize}

\subsubsection{4th-Generation Tensor Cores and FP8 Transformer Engine}
\label{sec:appendix_ada_tensor_cores}

The 4th-generation Tensor Cores in the AD104 architecture provide native microarchitectural support for 8-bit floating-point (FP8) numerical formats, managed through the hardware FP8 Transformer Engine:
\begin{enumerate}
    \item \textbf{Dual FP8 Numerical Representations}:
    \begin{itemize}
        \item \textbf{FP8 E4M3} (1 sign bit, 4 exponent bits, 3 mantissa bits; dynamic range $[-448, 448]$, precision $\epsilon = 0.125$): Designed for the forward pass activation and weight representations where higher numerical precision is required to maintain language model perplexity.
        \item \textbf{FP8 E5M2} (1 sign bit, 5 exponent bits, 2 mantissa bits; dynamic range extending from $-57{,}344$ to $57{,}344$, matching IEEE FP16): Designed for backward gradient propagation and optimizer states where wide dynamic range is critical to prevent underflow.
    \end{itemize}
    \item \textbf{Structural 2:4 Sparsity Hardware Engine}: The Tensor Core includes dedicated hardware logic for $2:4$ structured sparse matrix multiplication. In every contiguous block of four values along the matrix inner dimension, exactly two values must be zero. The hardware compresses the dense matrix by $50\%$ into a packed data array accompanied by a 2-bit metadata index per non-zero pair, doubling GEMM computational throughput to \textbf{$320.4\text{ TFLOPS}$ (FP16/BF16)} and \textbf{$640.8\text{ TFLOPS}$ (FP8)} on the RTX 4070.
    \item \textbf{Asynchronous Memory Transfer Engine (\texttt{cp.async}) and Warp Specialization}: While Hopper (SM 9.0a) introduces the hardware Tensor Memory Accelerator (TMA) to asynchronously stream 5D tensor boxes via hardware descriptors, Ada Lovelace relies on software-pipelined asynchronous copy instructions (\texttt{cp.async.ca} and \texttt{cp.async.cg}). These instructions bypass the thread register file entirely, transferring 16-byte aligned segments directly from global memory/L2 into shared memory:
    \begin{lstlisting}[language=C++]
    // Asynchronous 16-byte transfer from DRAM/L2 directly to shared memory
    asm volatile("cp.async.cg.shared.global [%0], [%1], 16;\n" 
                 :: "r"(smem_ptr), "l"(global_ptr));
    asm volatile("cp.async.commit_group;\n");
    asm volatile("cp.async.wait_group 0;\n");
    \end{lstlisting}
    This hardware bypass prevents register exhaustion while sustaining peak bandwidth. Under modern warp specialization paradigms, active thread block warps are structurally partitioned into dedicated \emph{producer warps} (which issue asynchronous \texttt{cp.async} transfers into double-buffered shared memory) and \emph{consumer warps} (which load operands from shared memory and execute matrix multiply-accumulate operations on Tensor Cores). This eliminates synchronization barrier stalls and hides the $200\text{--}400$ cycle DRAM access latency behind continuous Tensor Core compute.
\end{enumerate}

\subsubsection{Operational Roofline Model and Memory Bandwidth Asymmetry}
\label{sec:appendix_ada_roofline}

The memory subsystem of the RTX 4070 comprises a 192-bit bus interfacing six 32-bit channels of GDDR6X DRAM, yielding a theoretical peak bandwidth of \textbf{$504.0\text{ GB/s}$}. Empirical microbenchmarking via pointer-chasing and STREAM benchmarks measures a sustained global memory bandwidth of \textbf{$255.8\text{ GB/s}$}.

The severe asymmetry between the peak computing throughput of the ALUs and the external memory bus creates high arithmetic intensity ridge points (knee points) under the Williams Roofline Model~\cite{williams2009roofline}:
\begin{align}
    I^*_{\text{base}} &= \frac{\text{Peak FP32 Vector TFLOPS}}{\text{Peak Memory Bandwidth}} = \frac{70.8 \times 10^{12}\text{ FLOP/s}}{504.0 \times 10^9\text{ Bytes/s}} \approx \mathbf{140.5\text{ FLOP/Byte}}, \\
    I^* &= \frac{\text{Peak Tensor FP16 Boost TFLOPS}}{\text{Peak Memory Bandwidth}} = \frac{160.2 \times 10^{12}\text{ FLOP/s}}{504.0 \times 10^9\text{ Bytes/s}} \approx \mathbf{317.9\text{ FLOP/Byte}}, \\
    I^*_{\text{Tensor FP8}} &= \frac{\text{Peak Tensor FP8 Boost TFLOPS}}{\text{Peak Memory Bandwidth}} = \frac{320.4 \times 10^{12}\text{ FLOP/s}}{504.0 \times 10^9\text{ Bytes/s}} \approx \mathbf{635.7\text{ FLOP/Byte}}.
\end{align}
Empirical hardware counter auditing reveals an audited operational ridge point of $I^*_{\text{audit}} = 183.1\text{ FLOP/Byte}$. Because autoregressive linear attention decoding exhibits an arithmetic intensity between $2$ and $16\text{ FLOP/Byte}$, naive kernels operate deep within the memory-bandwidth-bound plateau. Achieving operational acceleration on AD104 silicon requires fusing sequence sweeps inside the $48\text{ MB}$ L2 cache and maximizing shared memory reuse.

% ==============================================================================
% D.2 Huawei Ascend 910B DaVinci NPU Architecture Dissection
% ==============================================================================
\subsection{Huawei Ascend 910B DaVinci NPU Architecture Dissection}
\label{sec:appendix_ascend_dissection}

\subsubsection{AI Core Heterogeneous Decomposition: Cube, Vector, and Scalar Units}
\label{sec:appendix_ascend_ai_core}

The Huawei Ascend 910B processor is constructed upon the proprietary DaVinci v2 Neural Processing Unit (NPU) microarchitecture, managed by the Compute Architecture for Neural Networks (CANN 8.0) software ecosystem. Unlike the homogeneous SIMT multithreading model of NVIDIA GPUs, each Ascend AI Core is architecturally disaggregated into three specialized, physically decoupled processing engines, as illustrated in Figure~\ref{fig:ascend_aicore_dataflow}:
\begin{enumerate}
    \item \textbf{Cube Unit (Matrix Multiplication Engine)}: The primary compute workhorse of the AI Core. The Cube Unit is a 3D systolic tensor array designed strictly for dense matrix multiply-accumulate operations. In each clock cycle, a single Cube Unit computes a \textbf{$16 \times 16 \times 16$ matrix multiplication} on FP16 or BF16 operands:
    \begin{equation}
        \mat{C}_{16 \times 16} = \mat{A}_{16 \times 16} \times \mat{B}_{16 \times 16} + \mat{C}_{16 \times 16},
    \end{equation}
    executing \textbf{4,096 MAC operations (8,192 FLOPs) per cycle}. At a nominal operating clock of $1.05\text{ GHz}$, the Cube Unit delivers \textbf{$280.0\text{ TFLOPS}$} of FP16/BF16 dense matrix compute per chip. The Cube Unit operates exclusively on tensors formatted in proprietary 5D fractal layouts (\texttt{Fractal NZ}).
    \item \textbf{Vector Unit (SIMD Processing Engine)}: An ultra-wide SIMD processor containing 128 32-bit vector ALUs and 256 16-bit vector ALUs. The Vector Unit delivers \textbf{$14.0\text{ TFLOPS}$ of FP32} and \textbf{$28.0\text{ TFLOPS}$ of FP16} throughput. It is responsible for all elementwise arithmetic, activation functions ($\text{GeLU}$, $\text{SiLU}$, $\text{ELU}$), normalization layers ($\text{LayerNorm}$, $\text{RMSNorm}$), and associative prefix scans required by state-space models.
    \item \textbf{Scalar Unit (Control and Address Engine)}: A dedicated processor responsible for loop counter maintenance, control-flow branching, instruction issue, and pointer base-offset generation. Because the Scalar Unit operates sequentially, excessive pointer manipulation or non-vectorized striding in custom kernels leads to \textbf{Scalar Pipeline Saturation}, where both the Cube and Vector units stall awaiting instruction issue.
\end{enumerate}

% ------------------------------------------------------------------------------
% TikZ Figure D.2: Ascend 910B AI Core Dataflow Schematic
% ------------------------------------------------------------------------------
\begin{figure*}[t]
\centering
\resizebox{0.98\linewidth}{!}{%
\begin{tikzpicture}[
    scale=0.88, every node/.style={transform shape},
    corebox/.style={draw=red!80!black, fill=red!3, rounded corners=4pt, line width=1.2pt},
    compbox/.style={draw=red!70!black, fill=red!10, rounded corners=3pt, line width=0.9pt, font=\footnotesize\bfseries, align=center},
    bufbox/.style={draw=blue!70!black, fill=blue!8, rounded corners=2pt, line width=0.8pt, font=\footnotesize, align=center},
    mtebox/.style={draw=purple!80!black, fill=purple!10, rounded corners=2pt, line width=0.9pt, font=\footnotesize\bfseries, align=center},
    hbmbox/.style={draw=green!60!black, fill=green!10, rounded corners=3pt, line width=1.1pt, font=\small\bfseries, align=center},
    arr/.style={-{Stealth[length=5pt,width=4pt]}, line width=1pt, draw=black!75},
    darr/.style={{Stealth[length=5pt,width=4pt]}-{Stealth[length=5pt,width=4pt]}, line width=1.1pt, draw=black!75}
]

% Main AI Core Boundary
\draw[corebox] (-0.5, 2.0) rectangle (15.5, 10.6);
\node[anchor=north west, font=\normalsize\bfseries, text=red!80!black] at (-0.3, 10.45) {Huawei Ascend 910B DaVinci v2 AI Core Subsystem (30--32 Cores)};

% Global Memory (HBM2e)
\draw[hbmbox] (1.0, 0.2) rectangle (14.0, 1.3);
\node at (7.5, 0.75) {High Bandwidth Memory (32 GB / 64 GB HBM2e on 4096-bit Bus, 1.2 -- 1.6 TB/s)};

% MTE2 and MTE3 DMA Engines
\node[mtebox, text width=5.5cm] (mte2) at (3.5, 2.8) {MTE2 (DMA Inbound Engine)\\{\tiny Global Memory (HBM) $\to$ UB / L1 Buffer}};
\node[mtebox, text width=5.5cm] (mte3) at (11.5, 2.8) {MTE3 (DMA Outbound Engine)\\{\tiny Unified Buffer (UB) $\to$ Global Memory (HBM)}};

% Arrows from HBM to MTE
\draw[arr] (3.5, 1.3) -- (mte2);
\draw[arr] (mte3) -- (11.5, 1.3);

% L1 Buffer
\node[bufbox, text width=4.0cm] (l1) at (3.5, 4.4) {L1 Matrix Buffer\\{\tiny 512 KB -- 1 MB SRAM (Weights/Cache)}};
\draw[arr] (mte2) -- (l1);

% L0A and L0B Buffers
\node[bufbox, text width=2.4cm] (l0a) at (1.8, 6.0) {L0A Buffer\\{\tiny 64 KB (Matrix A)}};
\node[bufbox, text width=2.4cm] (l0b) at (5.2, 6.0) {L0B Buffer\\{\tiny 64 KB (Matrix B)}};
\draw[arr] (l1) -- (l0a);
\draw[arr] (l1) -- (l0b);

% Cube Unit
\node[compbox, fill=red!18, text width=5.0cm] (cube) at (3.5, 7.8) {Cube Unit (3D Matrix MAC)\\{\tiny $16 \times 16 \times 16$ FP16 MACs (4,096 MACs/Cycle)\\{\scriptsize Peak: 280 TFLOPS (Fractal NZ Format)}}};
\draw[arr] (l0a) -- (cube);
\draw[arr] (l0b) -- (cube);

% L0C Accumulator
\node[bufbox, text width=4.0cm] (l0c) at (3.5, 9.6) {L0C Accumulator Buffer\\{\tiny 128 KB FP32 Accumulation}};
\draw[arr] (cube) -- (l0c);

% Unified Buffer (UB)
\node[bufbox, fill=blue!15, text width=5.0cm] (ub) at (11.5, 6.5) {\textbf{Unified Buffer (UB) Scratchpad}\\{\tiny 256 KB Local Fast SRAM (Up to 512 KB)\\{\scriptsize Software-Managed; Strict Physical Addressing}}};
\draw[arr] (mte2) -- (ub);
\draw[arr] (ub) -- (mte3);
\draw[arr] (l0c) -| node[above, font=\tiny\bfseries, text=black!70] {DataCopy} (ub);

% Vector Unit
\node[compbox, fill=orange!15, text width=4.8cm] (vec) at (11.5, 9.2) {Vector Unit (SIMD ALUs)\\{\tiny 128 $\times$ FP32 / 256 $\times$ FP16 ALUs\\{\scriptsize Elementwise, Reductions, Softmax, Scans}}};
\draw[darr] (ub) -- node[right, font=\tiny\bfseries] {Vector Load/Store} (vec);

% Scalar Unit
\node[compbox, fill=purple!12, text width=13.0cm] (scalar) at (7.5, 3.8) {Scalar Unit: Core Control, 64-bit Address Generation, Branch Dispatch, and Pipeline Queue Orchestration};

\end{tikzpicture}%
}
\caption[Huawei Ascend 910B (DaVinci v2) AI Core Pipeline and Memory Hierarchy]{\textbf{Huawei Ascend 910B (DaVinci v2) AI Core Pipeline and Memory Hierarchy.} The architecture decouples the matrix multiplication engine (Cube Unit) from elementwise and scan operations (Vector Unit). Data movement between off-chip HBM2e and on-chip SRAM buffers (L1, L0A/B, L0C, and the $256\text{ KB}$ Unified Buffer) is driven by asynchronous Memory Transfer Engines (MTE2 and MTE3) utilizing explicit double-buffered queue synchronization.}
\label{fig:ascend_aicore_dataflow}
\end{figure*}

\subsubsection{Internal Buffer Hierarchy and Physical Scratchpad Memory}
\label{sec:appendix_ascend_memory_hierarchy}

In sharp contrast to the unified, hardware-managed cache architecture of NVIDIA GPUs, the Ascend AI Core utilizes an explicit, physically partitioned scratchpad memory hierarchy:
\begin{enumerate}
    \item \textbf{Unified Buffer (UB)}: A $256\text{ KB}$ (expandable to $512\text{ KB}$ in DaVinci v2 Pro) high-speed on-chip SRAM array per AI Core. The UB serves as the universal staging area: all inputs and outputs for the Vector Unit, matrix results copied from the L0C accumulator, and intermediate data dispatched to or from global memory must reside in the UB. The UB has \textbf{no virtual memory abstraction} and \textbf{no hardware paging mechanism}. If the live tensor working set inside an operator exceeds $256\text{ KB}$, the hardware generates an unrecoverable \texttt{SIGBUS} / AIV memory fault, aborting execution immediately.
    \item \textbf{L1 Matrix Buffer}: A dedicated on-chip buffer of $512\text{ KB}$ to $1\text{ MB}$ utilized exclusively for staging weight tiles and feature maps before ingestion into the Cube Unit.
    \item \textbf{L0A and L0B Input Buffers}: Paired $64\text{ KB}$ buffers feeding the left ($\mat{A}$) and right ($\mat{B}$) operand matrices into the Cube Unit systolic array.
    \item \textbf{L0C Accumulator Buffer}: A $128\text{ KB}$ accumulator buffer holding the output matrix ($\mat{C}$) in FP32 precision. Data computed in L0C cannot be directly processed by the Vector Unit; it must first be explicitly copied to the UB via the internal bus using the \texttt{DataCopy} primitive.
    \item \textbf{Global Memory Subsystem}: 32 GB or 64 GB of High Bandwidth Memory (HBM2e) interfaced through four 1024-bit physical stacks (4096-bit composite bus), delivering \textbf{$1.2\text{ to }1.6\text{ TB/s}$ of sustained global memory bandwidth}.
\end{enumerate}

\subsubsection{Unified Buffer Capacity Planning and Live Footprint Budgeting}
\label{sec:appendix_ascend_ub_budgeting}

Because UB overflow results in catastrophic process termination, tile dimension budgeting is the fundamental correctness requirement for custom Ascend kernels. For a chunked recurrent linear attention kernel (e.g., \texttt{chunk\_kda}), the active memory footprint within the UB at any point during execution is governed by:
\begin{equation}
    \text{Footprint}_{\text{UB}} = \left[ 2 \cdot (B_T \cdot B_K + B_T \cdot B_V) + 2 \cdot (B_K \cdot B_V) + 2 \cdot B_T \right] \times \text{sizeof}(\text{dtype}) \times M_{\text{buf}},
    \label{eq:ascend_ub_budget_formula}
\end{equation}
where:
\begin{itemize}
    \item $B_T$ denotes the temporal sequence chunk tile size (e.g., $32, 64$).
    \item $B_K$ and $B_V$ denote the feature key and value tile dimensions ($64, 128$).
    \item $M_{\text{buf}} \in \{1, 2\}$ denotes the buffering multiplier ($M_{\text{buf}} = 2$ for double-buffered ping-pong pipelines).
    \item A mandatory \textbf{Safety Margin of 20\%} is reserved for compiler register allocation, stack spill headroom, and layout transformation buffers.
\end{itemize}

Under standard FP16 execution ($\text{sizeof}(\text{dtype}) = 2\text{ bytes}$), the usable capacity on a $256\text{ KB}$ UB core is:
\begin{equation}
    \text{Capacity}_{\text{usable}} = 256 \times 1,024 \times 0.80 = 209,715\text{ Bytes (204.8 KB)}.
\end{equation}
If a developer naively selects tile dimensions $B_T = 128, B_K = 128, B_V = 128$ under double buffering ($M_{\text{buf}} = 2$):
\begin{align}
    \text{Footprint} &= \left[ 2(128 \times 128 + 128 \times 128) + 2(128 \times 128) + 2(128) \right] \times 2 \times 2 \notag \\
    &= \left[ 2(32,768) + 32,768 + 256 \right] \times 4 = 98,560 \times 4 = \mathbf{394,240\text{ Bytes (385.0 KB)}}.
\end{align}
This allocation exceeds physical UB capacity by \textbf{$54.0\%$}, causing instant hardware failure. The \neuralops{} \texttt{AscendUBManager} dynamically solves this constrained integer program, enforcing tile constraints $B_T = 64, B_K = 64, B_V = 64$ ($104.5\text{ KB}$ live footprint), guaranteeing complete UB stability.

\subsubsection{Asynchronous DMA Double-Buffering Mechanics (MTE2/MTE3)}
\label{sec:appendix_ascend_dma_double_buffering}

Data movement between off-chip HBM and on-chip SRAM is orchestrated by two dedicated hardware engines:
\begin{itemize}
    \item \textbf{MTE2 (Memory Transfer Engine Inbound)}: Asynchronously reads contiguous or strided tensor blocks from Global Memory (HBM) and deposits them into the UB or L1 Buffer.
    \item \textbf{MTE3 (Memory Transfer Engine Outbound)}: Asynchronously flushes computed results from the UB back into Global Memory.
\end{itemize}

To completely hide memory access latency, \neuralops{} Ascend kernels partition the Unified Buffer into two symmetric ping-pong workspaces: \textbf{Buffer 0} and \textbf{Buffer 1}. Execution proceeds in a pipelined three-stage overlapping timeline, formalized in Table~\ref{tab:ascend_double_buffer_timeline}.

\begin{table}[t]
\centering
\small
\caption[Three-Stage Double-Buffered Pipeline Timeline on Ascend AI Core]{\textbf{Three-Stage Double-Buffered Pipeline Timeline on Ascend AI Core.} Memory movement via MTE2/MTE3 overlaps execution on Vector/Cube units across successive sequence tiles.}
\label{tab:ascend_double_buffer_timeline}
\setlength{\tabcolsep}{8pt}
\renewcommand{\arraystretch}{1.15}
\resizebox{\columnwidth}{!}{%
\begin{tabular}{cccc}
\toprule
\textbf{Pipeline Cycle} & \textbf{MTE2 Inbound DMA} & \textbf{Vector / Cube Compute} & \textbf{MTE3 Outbound DMA} \\
\midrule
Cycle 0 & Load Tile 0 $\to$ Buffer 0 & \emph{Stalled (Pipeline Warmup)} & \emph{Idle} \\
Cycle 1 & Load Tile 1 $\to$ Buffer 1 & Compute Tile 0 on Buffer 0 & \emph{Idle} \\
Cycle 2 & Load Tile 2 $\to$ Buffer 0 & Compute Tile 1 on Buffer 1 & Store Tile 0 $\to$ HBM \\
Cycle 3 & Load Tile 3 $\to$ Buffer 1 & Compute Tile 2 on Buffer 0 & Store Tile 1 $\to$ HBM \\
$\dots$ & $\dots$ & $\dots$ & $\dots$ \\
Drain Cycle & \emph{Idle (End of Sequence)} & Compute Tile $K$ on Buffer 1 & Store Tile $K-1$ $\to$ HBM \\
\bottomrule
\end{tabular}%
}
\end{table}

In Ascend C / CCE programming, this pipeline is enforced through the hardware queue primitives:
\begin{lstlisting}[language=C++]
// Double-buffered Queue Setup in Ascend C
TQue<QuePosition::VECIN, 2> inQueue;
TQue<QuePosition::VECOUT, 2> outQueue;

// Compute Stage
LocalTensor<half> inputTensor = inQueue.DeQue<half>();
LocalTensor<half> outputTensor = outQueue.AllocTensor<half>();
VectorAdd(outputTensor, inputTensor, ...);
outQueue.EnQue(outputTensor);
inQueue.FreeTensor(inputTensor);
\end{lstlisting}

\paragraph{Microarchitectural Grounding of the Triton Accumulator Identity Copy Invariant.} During the synthesis of chunked recurrent linear attention kernels (e.g., \texttt{chunk\_kda}) under the \texttt{triton\_ascend} backend, a critical memory clobbering anomaly emerges from the interaction between asynchronous MTE2 DMA transfers and the BiSheng compiler's liveness analysis. In recurrent state update operations of the form:
\begin{equation}
    \mat{S}_c = \mat{S}_{c-1} + \mat{K}_c^\top \mat{V}_c,
\end{equation}
the inner systolic loop evaluates the outer-product update via \texttt{tl.dot(tl.trans(k\_buf0), v\_buf0)}. 

Under the double-buffered execution model, while the Cube Unit executes the matrix multiply-accumulate on tile $c$ resident in buffer partition $\text{UB}_0$, the MTE2 DMA engine concurrently streams tile $c+1$ from off-chip HBM into partition $\text{UB}_1$. However, because the BiSheng compiler performs liveness analysis at the level of scalar LLVM IR value definitions without tracking hardware DMA transfer latency, it detects that \texttt{k\_buf0} has no downstream consumers following ingestion into \texttt{tl.dot}. The compiler's physical memory reuse pass consequently flags the SRAM region occupied by \texttt{k\_buf0} in $\text{UB}_0$ as dead and immediately reclaimable.

Crucially, the physical systolic array of the Cube Unit does not read operands directly from UB; instead, an internal sub-engine (MTE1) asynchronously transfers the matrix blocks from $\text{UB}_0$ into the L0A and L0B matrix staging buffers. Because this internal DMA transfer executes concurrently with both the Vector Unit and MTE2, a destructive hardware race condition occurs: the scheduler launches an inbound MTE2 transfer for the subsequent tile's key vector $\mat{K}_{c+1}$ into the address range marked as free, overwriting the trailing 32-byte blocks of $\mat{K}_c$ before MTE1 has completed transferring them into L0A. This memory clobbering produces silent numerical divergence, manifesting as $\text{NaN}$ or $\pm\infty$ values in the recurrent state accumulator.

To resolve this microarchitectural conflict, \neuralops{} establishes the mandatory \textbf{\texttt{+ 0.0} Identity Copy Invariant}:
\begin{lstlisting}[language=Python]
# Enforce explicit Vector ALU copy in UB to isolate physical memory from MTE2 DMA
k_compute = k_buf0 + 0.0
s_update = tl.dot(tl.trans(k_compute), v_buf0)
\end{lstlisting}
By introducing an elementwise addition with floating-point zero, the kernel forces the generation of an explicit Vector Unit instruction (\texttt{VADD}). In the SSA representation, this establishes a new live value range with an independent physical destination buffer in UB. The compiler is thereby prohibited from reclaiming the memory of \texttt{k\_buf0} until the Vector Unit instruction completes and synchronizes across the internal pipeline barrier (\texttt{pipe\_barrier(PIPE\_V)}), guaranteeing that $\mat{K}_c$ remains intact throughout its asynchronous staging into L0A.

% ==============================================================================
% D.3 Comparative Analysis: CUDA C++ / Triton vs. Ascend C / CANN
% ==============================================================================
\subsection{Comparative Analysis: CUDA C++ / Triton vs. Ascend C / CANN}
\label{sec:appendix_comparative_analysis}

\subsubsection{Memory Models and Spatial Alignment Contracts}
\label{sec:appendix_memory_models}

The divergent execution philosophies of NVIDIA and Huawei hardware impose distinct spatial alignment constraints on memory allocations:
\begin{itemize}
    \item \textbf{NVIDIA CUDA Alignment}: Enforces natural 4-byte and 16-byte alignment to enable vectorized 128-bit memory instructions (\texttt{LDG.E.128} / \texttt{STG.E.128} in SASS). Stride padding focuses on avoiding the 32-way shared memory bank collisions.
    \item \textbf{Huawei Ascend Alignment}: Imposes two mandatory hardware constraints:
    \begin{enumerate}
        \item \textbf{32-Byte Alignment Contract}: Every tensor base address, row stride, and DMA block transfer must be an exact multiple of \textbf{32 bytes} (equivalent to 16 FP16 elements or 8 FP32 elements). Memory reads violating 32-byte alignment trigger hardware alignment traps that abort execution.
        \item \textbf{512-Byte Burst Contract}: To attain $\ge 90\%$ of peak bus bandwidth across the HBM2e memory controller, DMA burst transfers executed by MTE2/MTE3 must align to \textbf{512-byte boundaries} (256 FP16 elements). Transfers below 512 bytes degrade bus efficiency by up to $8\times$.
    \end{enumerate}
    \item \textbf{Fractal NZ vs. Dense ND Layouts}: NVIDIA Tensor Cores process conventional 2D row-major or column-major matrices (ND format) directly through warp-level fragment loads (\texttt{wmma::load\_matrix\_sync}). Conversely, the Ascend Cube Unit cannot ingest ND matrices. All tensors must be transposed into 5D Fractal NZ format:
    \begin{equation}
        \text{Shape}_{\text{ND}} = [M, K] \quad \xrightarrow{\text{Fractal Transformation}} \quad \text{Shape}_{\text{NZ}} = \left[ \left\lceil \frac{M}{16} \right\rceil, \left\lceil \frac{K}{16} \right\rceil, 16, 16 \right].
    \end{equation}
    In custom kernels, failing to select tile dimensions as exact multiples of 16 forces the CANN runtime to inject dynamic layout conversion kernels (\texttt{TransData}), completely eroding execution speedups.
\end{itemize}

\subsubsection{Compilation Toolchains and Intermediate Representations}
\label{sec:appendix_compilation_toolchains}

The software compilation toolchains reflect the underlying execution divergence:
\begin{itemize}
    \item \textbf{NVIDIA Execution Stack}: The source CUDA C++ code is parsed by \texttt{nvcc}, generating Parallel Thread Execution (PTX) virtual assembly, which is subsequently compiled by the \texttt{ptxas} optimizing assembler into machine-level SASS instructions. In the Level 2 Triton tier, Python ASTs are compiled into Triton-IR, lowered through LLVM-IR to NVVM, and emitted as PTX. Dynamic register allocation and warp scheduling are resolved at compile time with fine-grained control via \texttt{\_\_launch\_bounds\_\_}.
    \item \textbf{Huawei Ascend Stack}: The Ascend C / CANN ecosystem employs the \textbf{BiSheng Compiler} (an LLVM-derived proprietary compiler targeting the Compute Core Engine, CCE). Triton kernels compiled via the \texttt{triton\_ascend} backend bypass PTX entirely, compiling Triton-IR directly into CCE C++ extensions. These extensions explicitly manage hardware event synchronization flags (\texttt{SetFlag} / \texttt{WaitFlag}) to enforce synchronization across the Scalar, Vector, Cube, and MTE units.
\end{itemize}

\subsubsection{Portability Bottlenecks: Non-Power-of-Two Head Dimensions}
\label{sec:appendix_portability_bottlenecks}

A critical architectural incompatibility between the two platforms emerges when deploying attention mechanisms with non-power-of-two head dimensions ($d_k \notin \{32, 64, 128, 256\}$), such as $d_k = 48, 80,$ or $96$ encountered in pruned models:
\begin{itemize}
    \item \textbf{On NVIDIA AD104}: Thread blocks handle arbitrary $d_k$ dimensions via simple thread-predicated masks (\texttt{threadIdx.x < d\_k}). Shared memory padding is dynamically adjusted ($\text{PAD} = 32 - (d_k \bmod 32)$) to maintain bank conflict immunity. Performance scales smoothly with the reduction in arithmetic operations.
    \item \textbf{On Huawei Ascend 910B}: If $d_k$ is not a multiple of 16, the Cube Unit cannot construct full $16 \times 16$ fractal blocks. The CANN compiler must either zero-pad the tensors to the next multiple of 16 (incurring up to $33\%$ memory and compute overhead) or fallback to executing the entire matrix multiplication on the **Vector Unit**. Because the Vector Unit peak throughput ($28\text{ TFLOPS}$) is exactly **$10\times$ slower** than the Cube Unit ($280\text{ TFLOPS}$), non-power-of-two head projections suffer severe performance degradation.
\end{itemize}

% ==============================================================================
% D.4 Master Microarchitectural Comparison Table
% ==============================================================================
\subsection{Master Microarchitectural Comparison Table}
\label{sec:appendix_master_comparison_table}

Table~\ref{tab:microarchitecture_master_dissection} provides a comprehensive, side-by-side microarchitectural comparison across the primary compute substrates evaluated in \neuralops{}, including datacenter reference architectures (NVIDIA A100 SXM4 and H100 SXM5).

\begin{table*}[t]
\centering
\scriptsize
\caption[Comprehensive Microarchitectural and Silicon Dissection]{\textbf{Comprehensive Microarchitectural and Silicon Dissection: NVIDIA AD104 vs. Huawei Ascend 910B with Datacenter Reference Baselines (Ampere A100 and Hopper H100).}}
\label{tab:microarchitecture_master_dissection}
\setlength{\tabcolsep}{3.0pt}
\renewcommand{\arraystretch}{1.15}
\resizebox{\textwidth}{!}{%
\begin{tabular}{lcccc}
\toprule
\textbf{Architectural Parameter} &
\textbf{GeForce RTX 4070 (AD104)} &
\textbf{Ascend 910B (DaVinci v2)} &
\textbf{A100 SXM4 (Ampere)} &
\textbf{H100 SXM5 (Hopper)} \\
\midrule
Compute Architecture & NVIDIA Ada Lovelace (SM 8.9) & Huawei DaVinci v2 & NVIDIA Ampere (SM 8.0) & NVIDIA Hopper (SM 9.0a) \\
Process Technology Node & TSMC 4N ($5\text{ nm}$ Custom) & $7\text{ nm}$ FinFET Process & TSMC $7\text{ nm}$ & TSMC 4N ($5\text{ nm}$ Custom) \\
Transistor Count & 35.8 Billion & $\approx 32.0\text{ Billion}$ & 54.2 Billion & 80.0 Billion \\
Physical Die Size & $294.5\text{ mm}^2$ & $\approx 400.0\text{ mm}^2$ & $826.0\text{ mm}^2$ & $814.0\text{ mm}^2$ \\
Total Processing Clusters & 46 SMs (60 SMs in Ti) & 30--32 AI Cores & 108 SMs & 132 SMs \\
Base / Boost Clock & $1,695\text{ / }2,175\text{ MHz}$ & $1,000\text{ / }1,050\text{ MHz}$ & $1,095\text{ / }1,410\text{ MHz}$ & $1,590\text{ / }1,980\text{ MHz}$ \\
\midrule
\textbf{Compute Core Composition} & & & & \\
--- FP32 Standard ALUs & 5,888 CUDA Cores & 32 Vector Cores (SIMD) & 6,912 CUDA Cores & 16,896 CUDA Cores \\
--- Specialized Matrix Cores & 184 4th-Gen Tensor Cores & 32 Cube Units ($16^3$ MAC) & 432 3rd-Gen Tensor Cores & 528 4th-Gen Tensor Cores \\
--- Dual-Issue Pipeline & Concurrent FP32 + INT32 & Disaggregated Cube/Vector & Independent FP32/INT32 & Concurrent FP32 + INT32 \\
\midrule
\textbf{Peak Compute Throughput} & & & & \\
--- FP32 Non-Tensor & $70.8\text{ TFLOPS}$ (Vector Limit) & $32.0\text{ TFLOPS}$ (Cube+Vec) & $19.5\text{ TFLOPS}$ & $67.0\text{ TFLOPS}$ \\
--- FP16 / BF16 Dense Tensor & $160.2\text{ TFLOPS}$ (Boost) / $141.6\text{ (Base)}$ & $280.0\text{ TFLOPS}$ & $312.0\text{ TFLOPS}$ & $989.0\text{ TFLOPS}$ \\
--- FP16 / BF16 2:4 Sparse & $320.4\text{ TFLOPS}$ & Not Supported & $624.0\text{ TFLOPS}$ & $1,978.0\text{ TFLOPS}$ \\
--- FP8 Native Tensor & $320.4\text{ TFLOPS}$ (Dense) / $640.8\text{ (Sparse)}$ & Quantized Simulation Only & Not Supported & $1,978.0\text{ TFLOPS}$ \\
\midrule
\textbf{On-Chip SRAM Hierarchy} & & & & \\
--- Register File per Core & 64K 32-bit ($256\text{ KB}$) & Integrated Vector Registers & 64K 32-bit ($256\text{ KB}$) & 64K 32-bit ($256\text{ KB}$) \\
--- Total Chip Register Capacity & $11.78\text{ MB}$ ($15.36\text{ MB}$ full) & $\approx 8.0\text{ MB}$ & $27.65\text{ MB}$ & $33.79\text{ MB}$ \\
--- Scratchpad / Shared Memory & 128 KB (Max 99 KB Smem) & $256\text{--}512\text{ KB}$ (UB) & 192 KB (Max 164 KB Smem) & 256 KB (Max 228 KB Smem) \\
--- Matrix Buffers & None (Unified SRAM) & 1 MB L1 + 256 KB L0A/B/C & None (Unified SRAM) & 228 KB Shared + TMA \\
--- On-Chip L2 Cache & \textbf{36 MB} (48 MB in Ti) & Distributed Cross-Core Cache & 40 MB & 50 MB \\
--- L2 Cache Bandwidth & $\ge 1.8\text{ TB/s}$ & $\approx 1.2\text{ TB/s}$ & $\approx 2.4\text{ TB/s}$ & $\approx 4.8\text{ TB/s}$ \\
\midrule
\textbf{Off-Chip Memory Subsystem} & & & & \\
--- Memory Technology & GDDR6X (Micron PAM4) & HBM2e (4 Stacks) & HBM2e (5 Stacks) & HBM3 (6 Stacks) \\
--- Memory Capacity & 12 GB & 32 GB / 64 GB & 80 GB & 80 GB \\
--- Memory Bus Interface & 192-bit & 4096-bit & 5120-bit & 5120-bit \\
--- Theoretical Bandwidth & $504.0\text{ GB/s}$ & $1,200\text{--}1,600\text{ GB/s}$ & $2,039.0\text{ GB/s}$ & $3,350.0\text{ GB/s}$ \\
--- Audited STREAM Bandwidth & $255.8\text{ GB/s}$ & $1,120.0\text{ GB/s}$ & $1,780.0\text{ GB/s}$ & $2,920.0\text{ GB/s}$ \\
\midrule
\textbf{Asynchronous Transfer Engine} & \texttt{cp.async} (Software Pipeline) & MTE2 (In) / MTE3 (Out) DMA & \texttt{cp.async} (Gen 1) & Hardware TMA + DSMEM \\
\textbf{Hardware Tensor Layout} & Dense ND / 2:4 Sparse & 5D Fractal NZ ($16 \times 16$) & Dense ND / 2:4 Sparse & Dense ND / Swizzled \\
\textbf{Arithmetic Intensity Ridge ($I^*$)} & $140.5\text{ (FP32) / }317.9\text{ (FP16 Boost)}$ & $\approx 140.0\text{ FLOP/Byte}$ & $153.0\text{ FLOP/Byte}$ & $295.2\text{ FLOP/Byte}$ \\
\textbf{Software Stack} & CUDA 12.8 / Triton 3.6 & CANN 8.0 / torch\_npu & CUDA 12.2 / Triton 3.0 & CUDA 12.6 / Triton 3.4 \\
\bottomrule
\end{tabular}%
}
\end{table*}

Table~\ref{tab:on_chip_sram_limits} summarizes the operational boundaries, allocation limits, and clock latencies across the memory hierarchy levels of both architectures.

\begin{table}[t]
\centering
\small
\caption[On-Chip SRAM and Buffer Allocation Limits]{\textbf{On-Chip SRAM and Buffer Allocation Limits, Latencies, and Bandwidths Across Memory Levels.}}
\label{tab:on_chip_sram_limits}
\setlength{\tabcolsep}{5pt}
\renewcommand{\arraystretch}{1.12}
\resizebox{\columnwidth}{!}{%
\begin{tabular}{lcccc}
\toprule
\textbf{Memory Level} & \textbf{Silicon Target} & \textbf{Size per SM / Core} & \textbf{Peak Bandwidth} & \textbf{Access Latency} \\
\midrule
Register File & AD104 (Ada) & 256 KB (64K $\times$ 32-bit) & $\approx 900\text{ GB/s per SM}$ & 1 clock cycle \\
Shared Memory & AD104 (Ada) & Configurable up to 99 KB & $128\text{ Bytes/cycle per SM}$ & 19--28 clock cycles \\
L1 Data Cache & AD104 (Ada) & Configurable (29--128 KB) & Shared with Smem & 28--34 clock cycles \\
L2 Crossbar Cache & AD104 (Ada) & 36 MB / 48 MB (Chip) & $\ge 1.8\text{ TB/s (Composite)}$ & 180--220 clock cycles \\
GDDR6X DRAM & AD104 (Ada) & 12 GB (Off-chip) & $504.0\text{ GB/s (Bus)}$ & 400--600 clock cycles \\
\midrule
Unified Buffer (UB) & Ascend 910B & 256 KB -- 512 KB & $\approx 800\text{ GB/s per Core}$ & 3--5 clock cycles \\
L1 Matrix Buffer & Ascend 910B & 512 KB -- 1 MB & $\approx 600\text{ GB/s per Core}$ & 12--16 clock cycles \\
L0A / L0B Buffers & Ascend 910B & 64 KB each (Matrix Inputs) & Direct to Systolic Array & 1 clock cycle \\
L0C Accumulator & Ascend 910B & 128 KB (FP32 Accum) & Systolic Feedback & 1 clock cycle \\
HBM2e Memory & Ascend 910B & 32 GB / 64 GB (Off-chip) & $1,200\text{--}1,600\text{ GB/s}$ & 250--350 clock cycles \\
\bottomrule
\end{tabular}%
}
\end{table}

% ==============================================================================
% D.5 Empirical Telemetry & Architectural Case Studies
% ==============================================================================
\subsection{Empirical Telemetry and Architectural Case Studies on RTX 4070}
\label{sec:appendix_empirical_telemetry}

To demonstrate how the physical silicon constraints established in Sections~\ref{sec:appendix_nvidia_ada_dissection} and~\ref{sec:appendix_ascend_dissection} govern real-world performance, Table~\ref{tab:master_18_ops_telemetry} presents the complete empirical telemetry for \textbf{18 computational operators} across the 3-level fallback hierarchy of \neuralops{} on the NVIDIA GeForce RTX 4070 (SM 8.9). All benchmarks were executed under standardized evaluation dimensions ($B=2, H=12, D=768, N=512$, FP16 precision, CUDA 12.8) across 100 cold-cache timed trials. The complete benchmark ledger, raw performance logs, and NCU hardware traces are permanently cataloged in the open-access Hugging Face repository at \url{https://huggingface.co/datasets/Hooshaai/neuralops-results}.

\begin{table*}[t]
\centering
\tiny
\caption[Empirical Telemetry Across Three-Level Fallback Stack]{\textbf{Comprehensive Empirical Telemetry Across the 3-Level Fallback Stack on NVIDIA GeForce RTX 4070 (AD104, SM 8.9).} Evaluated at $B=2, H=12, D=768, N=512$, FP16 precision, CUDA 12.8. Telemetry captures median wall-clock execution latency across 100 trials, relative speedups, on-chip SRAM allocations, register allocations per thread, launch grid bounds, and primary microarchitectural bottlenecks. Complete multi-run traces and hardware telemetry are openly accessible at \url{https://huggingface.co/datasets/Hooshaai/neuralops-results}.}
\label{tab:master_18_ops_telemetry}
\setlength{\tabcolsep}{2.2pt}
\renewcommand{\arraystretch}{1.15}
\resizebox{\textwidth}{!}{%
\begin{tabular}{rlrrrrrrrcccp{4.2cm}}
\toprule
\textbf{\#} & \textbf{Kernel Operator} &
\textbf{PyTorch (ms)} & \textbf{Triton (ms)} & \textbf{CUDA (ms)} &
\textbf{Speedup (Cu/Py)} & \textbf{Speedup (Cu/Tri)} &
\textbf{Py VRAM} & \textbf{Cu VRAM} &
\textbf{SRAM} & \textbf{Regs} & \textbf{Launch Grid} & \textbf{Primary Microarchitectural Bottleneck} \\
\midrule
1  & \texttt{moe\_routing}                & $29.596$ & ---      & $0.022$ & $\mathbf{1340.6\times}$ & ---          & $44.05\text{ MB}$ & $44.05\text{ MB}$ & $4.0\text{ KB}$  & 32 & \texttt{Blk(128)}  & Register-resident; warp butterfly shuffle reduction \\
2  & \texttt{hgrn2\_fused\_bwd\_scan}      & $5.695$  & $0.031$  & $0.013$ & $\mathbf{435.5\times}$  & $2.35\times$ & $40.91\text{ MB}$ & $39.40\text{ MB}$ & $16.0\text{ KB}$ & 40 & \texttt{Blk(128)}  & SRAM fused; backward norm and gate pipelining \\
3  & \texttt{linear\_attention}            & $41.680$ & $0.266$  & $0.263$ & $\mathbf{158.7\times}$  & $1.01\times$ & $43.87\text{ MB}$ & $43.86\text{ MB}$ & $32.0\text{ KB}$ & 56 & \texttt{Blk(128)}  & L2 cache bound; intra-chip cumulative state scan \\
4  & \texttt{selective\_scan}             & $27.095$ & $13.648$ & $0.273$ & $\mathbf{99.1\times}$   & $\mathbf{49.9\times}$ & $10.43\text{ MB}$ & $10.91\text{ MB}$ & $48.0\text{ KB}$ & 64 & \texttt{Blk(128)}  & Dual-buffer shared memory; \texttt{(128,2)} bounds \\
5  & \texttt{chunked\_delta\_gate\_scan}   & $43.506$ & $0.127$  & $0.466$ & $93.4\times$            & $0.27\times$ & $38.79\text{ MB}$ & $38.78\text{ MB}$ & $32.0\text{ KB}$ & 60 & \texttt{Blk(64,64)}& Triton autotuned 2D tile scheduling superior \\
6  & \texttt{paged\_attention}            & $1.654$  & ---      & $0.1625$& $\mathbf{10.18\times}$  & ---          & $36.81\text{ MB}$ & $36.05\text{ MB}$ & $8.0\text{ KB}$  & 36 & \texttt{Blk(128)}  & Shared memory lookup table for KV page indexing \\
7  & \texttt{sliding\_window\_attn}        & $1.822$  & $0.252$  & $0.251$ & $7.3\times$             & $1.00\times$ & $119.65\text{ MB}$& $26.90\text{ MB}$ & $32.0\text{ KB}$ & 48 & \texttt{Blk(256)}  & $77.5\%$ VRAM reduction via SRAM band tiling \\
8  & \texttt{lut\_codebook\_dequant\_gemm} & $0.084$  & $0.025$  & $0.019$ & $4.38\times$            & $1.31\times$ & $44.05\text{ MB}$ & $43.39\text{ MB}$ & $16.0\text{ KB}$ & 32 & \texttt{Blk(256)}  & SRAM codebook staging; dequantization on-the-fly \\
9  & \texttt{block\_sparse\_attn}          & $1.816$  & $1.526$  & $0.364$ & $5.0\times$             & $4.19\times$ & $119.15\text{ MB}$& $26.90\text{ MB}$ & $24.0\text{ KB}$ & 44 & \texttt{Blk(128)}  & Eliminates zero blocks; $77.4\%$ VRAM memory save \\
10 & \texttt{fused\_rotary\_emb}           & $0.028$  & ---      & $0.008$ & $3.43\times$            & ---          & $47.22\text{ MB}$ & $45.22\text{ MB}$ & $4.0\text{ KB}$  & 28 & \texttt{Blk(128)}  & Vectorized \texttt{half2} math; zero intermediate VRAM \\
11 & \texttt{flash\_gqa}                   & $0.192$  & $0.032$  & $0.079$ & $2.43\times$            & $0.41\times$ & $60.09\text{ MB}$ & $34.02\text{ MB}$ & $48.0\text{ KB}$ & 52 & \texttt{Blk(128)}  & KV head sharing; Triton register reuse advantage \\
12 & \texttt{kronecker\_matmul}            & $0.030$  & $0.026$  & $0.017$ & $1.79\times$            & $1.51\times$ & $43.53\text{ MB}$ & $43.34\text{ MB}$ & $8.0\text{ KB}$  & 32 & \texttt{Blk(256)}  & Separable Kronecker factorization in shared memory \\
13 & \texttt{cosformer\_linear\_attention} & $0.118$  & $0.027$  & $0.106$ & $1.12\times$            & $0.26\times$ & $41.78\text{ MB}$ & $39.65\text{ MB}$ & $24.0\text{ KB}$ & 48 & \texttt{Blk(128)}  & SFU trigonometric mapping; Triton block unroll \\
\midrule
14 & \texttt{fused\_lowrank\_gated\_gemm}  & $0.027$  & $0.041$  & $0.870$ & $0.031\times$           & $0.047\times$ & $38.40\text{ MB}$ & $38.34\text{ MB}$ & $16.0\text{ KB}$ & 40 & \texttt{Blk(128)}  & Un-autotuned CUDA fallback ($0.870\,\text{ms}$); Triton JIT optimal ($0.041\,\text{ms}$) \\
15 & \texttt{w4a16\_gemm}                  & $0.104$  & ---      & $0.266$ & $0.39\times$            & ---          & $48.02\text{ MB}$ & $24.52\text{ MB}$ & $32.0\text{ KB}$ & 64 & \texttt{Blk(256)}  & $51\%$ VRAM reduction; ALU bit-unpacking latency \\
16 & \texttt{fused\_cross\_entropy}        & $0.012$  & $0.065$  & $0.037$ & $0.31\times$            & $1.75\times$ & $47.74\text{ MB}$ & $45.75\text{ MB}$ & $16.0\text{ KB}$ & 36 & \texttt{Blk(256)}  & PyTorch cuBLAS hand-tuned assembly outperforms \\
17 & \texttt{fp8\_gemm}                    & $0.078$  & $0.068$  & $0.507$ & $0.15\times$            & $0.13\times$ & $31.02\text{ MB}$ & $30.02\text{ MB}$ & $64.0\text{ KB}$ & 80 & \texttt{Blk(128)}  & \textbf{Register pressure (80 regs); 50\% occupancy cliff} \\
18 & \texttt{fused\_norm\_linear}          & $0.015$  & $0.019$  & $0.195$ & $0.08\times$            & $0.10\times$ & $41.34\text{ MB}$ & $41.34\text{ MB}$ & $8.0\text{ KB}$  & 48 & \texttt{Blk(128)}  & Strided non-coalesced reduction memory access \\
\bottomrule
\end{tabular}%
}
\end{table*}

\subsubsection{Case Study 1: The \texttt{fp8\_gemm} Occupancy Cliff Anomaly}
\label{sec:appendix_case_study_fp8}

A prominent empirical finding in Table~\ref{tab:master_18_ops_telemetry} is the substantial slowdown of the custom \texttt{fp8\_gemm} kernel, which achieved a relative speedup of only \textbf{$0.15\times$} compared to PyTorch (i.e., $6.7\times$ slower than the baseline), despite the theoretical peak throughput of $373\text{ TFLOPS}$ provided by the 4th-generation Tensor Cores.

A low-level microarchitectural autopsy reveals the physical cause:
\begin{enumerate}
    \item \textbf{Dynamic Dequantization Overhead}: The custom Level 1 kernel implements per-tensor dynamic scaling ($s_A, s_B$) to convert FP8 values into intermediate precision prior to accumulation. In CUDA SASS assembly, byte extraction (\texttt{\_\_byte\_perm}), bit-shifting, and IEEE FP32 float conversion require $18$ to $24$ scalar auxiliary registers per thread.
    \item \textbf{Occupancy Degradation}: Register allocation climbs to \textbf{80 registers per thread}. Referring to Table~\ref{tab:register_occupancy_breakdown}, an allocation of 80 registers per thread restricts each SM to a maximum of \textbf{24 active warps}, dropping theoretical SM occupancy to exactly \textbf{$50.0\%$}.
    \item \textbf{Memory Latency Unhiding}: At $50\%$ occupancy, the warp scheduler lacks an adequate pool of active warps to issue instructions during memory stalls. Because GDDR6X memory accesses take 200--400 cycles, the execution ALUs remain idle, spending over $68.4\%$ of clock cycles in stall states.
    \item \textbf{Comparison with PyTorch \texttt{torch.\_scaled\_mm}}: The PyTorch baseline dispatches to vendor-optimized cuBLASLt GEMM routines hand-written in machine assembly. These micro-kernels constrain register consumption to $36$ registers per thread, sustaining $\ge 75\%$ occupancy and utilizing multi-stage asynchronous copy pipelines to completely hide global memory latency.
\end{enumerate}

\subsubsection{Case Study 2: SVD Low-Rank Register Relief in Recurrent State Tracking}
\label{sec:appendix_case_study_svd_relief}

In recurrent state-space models and linear attention mechanisms (\texttt{RetNet}, \texttt{GLA}, \texttt{selective\_scan}), maintaining the recurrent hidden state tensor:
\begin{equation}
    \mat{S}_t \in \R^{d \times d}, \quad \text{where } d = 128,
\end{equation}
creates a physical bottleneck for the GPU register file. The full state matrix contains $128 \times 128 = 16,384$ numerical elements ($64\text{ KB}$ per head in FP32).

If an algorithm attempts to hold this state across a single 32-thread warp, each thread must store:
\begin{equation}
    \frac{16,384\text{ elements}}{32\text{ threads}} = \mathbf{512\text{ registers per thread}}.
\end{equation}
Because the physical hardware architecture enforces an absolute ceiling of \textbf{255 registers per thread}, the compiler is forced to spill $257+$ registers per thread into off-chip Local Memory, degrading recurrent execution throughput by up to two orders of magnitude.

\neuralops{} resolves this bottleneck by restructuring the recurrent state via \textbf{Singular Value Decomposition (SVD) Low-Rank Factorization}:
\begin{equation}
    \mat{S}_t \approx \mat{U}_t \bm{\Sigma}_t \mat{V}_t^\top, \quad \mat{U}_t \in \R^{d \times r}, \quad \mat{V}_t \in \R^{d \times r}, \quad \bm{\Sigma}_t \in \R^{r \times r}, \quad \text{with } r = 16 \ll d.
\end{equation}
The microarchitectural impact of this mathematical reformulation is immediate:
\begin{enumerate}
    \item The tall projection bases $\mat{U}_t$ and $\mat{V}_t$ are tiled in on-chip \textbf{Shared Memory (SRAM)}. Their combined footprint is:
    \begin{equation}
        2 \times (128 \times 16) \times 2\text{ bytes} = \mathbf{8.19\text{ KB}},
    \end{equation}
    which easily fits within the $99\text{ KB}$ shared memory budget of the SM.
    \item The active recurrence loop executed by the warp ALUs updates only the inner singular value core $\bm{\Sigma}_t \in \R^{16 \times 16}$, which contains only $256$ elements ($1,024\text{ bytes}$ in FP32).
    \item Distributing 256 elements across 32 warp threads requires exactly:
    \begin{equation}
        \frac{256\text{ elements}}{32\text{ threads}} = \mathbf{8\text{ scalar accumulator registers per thread}}.
    \end{equation}
    \item Accounting for loop index variables, pointer base addresses, and temporary math operands, total compiler register allocation drops from an impossible $512$ registers to \textbf{32 registers per thread}.
\end{enumerate}
This yields a \textbf{$16\times$ register relief}, completely eliminating local memory spilling ($0\text{ bytes spilled}$) and restoring SM occupancy to the theoretical ceiling of \textbf{$100.0\%$ (48 active warps)}. This register relief explains the substantial speedups observed in Table~\ref{tab:master_18_ops_telemetry} for recurrent operators, including $\mathbf{99.1\times}$ for \texttt{selective\_scan} and $\mathbf{435.5\times}$ for \texttt{hgrn2\_fused\_bwd\_scan}.

\subsubsection{Case Study 3: VRAM Materialization Elimination via SRAM Tiling}
\label{sec:appendix_case_study_vram_savings}

Beyond computational speedups, Table~\ref{tab:master_18_ops_telemetry} highlights significant reductions in peak global memory footprint:
\begin{itemize}
    \item In \texttt{sliding\_window\_attn} and \texttt{block\_sparse\_attn}, standard PyTorch implementations materialize the complete $B \times H \times N \times N$ attention affinity matrix in DRAM, populating elements outside the band with $-\infty$ before executing softmax. This requires \textbf{$119.65\text{ MB}$ of VRAM}.
    \item The \neuralops{} Level 1 CUDA kernel streams $128 \times 128$ tiles into SRAM, computes intermediate softmax exponentials on-the-fly, and commits only final output activations to DRAM, consuming only \textbf{$26.90\text{ MB}$ of VRAM}---a \textbf{$77.5\%$ memory saving}.
    \item In \texttt{w4a16\_gemm}, packed 4-bit integer weights are streamed directly into shared memory and unpacked into FP16 formats within ALU registers, reducing model parameter footprint from $48.02\text{ MB}$ to $24.52\text{ MB}$ (\textbf{$51.0\%$ reduction}).
\end{itemize}

\subsubsection{Case Study 4: Compiler Autotuning Disparity in \texttt{fused\_lowrank\_gated\_gemm}}
\label{sec:appendix_case_study_fused_lowrank}

Table~\ref{tab:master_18_ops_telemetry} reveals a profound microarchitectural dichotomy in operator fusion: \texttt{fused\_lowrank\_gated\_gemm} achieves its optimal execution under the Level 2 Triton JIT backend ($0.041\,\text{ms}$), delivering a $1.19\times$ speedup over the standalone PyTorch baseline ($0.051\,\text{ms}$) in isolated microbenchmarks. Conversely, the initial un-autotuned Level 1 CUDA C++ fallback kernel incurs significant grid scheduling and thread coarsening overhead on thin bottleneck matrices, executing in $0.870\,\text{ms}$ (a $0.0314\times$ speedup relative to dense PyTorch, documented as record \texttt{6975a7fc689cc386} in the master ledger). 

This empirical disparity underscores the design rationale of the \neuralops{} 3-level fallback architecture: compiler-autotuned Triton JIT kernels dynamically discover optimal block shapes (\texttt{BLOCK\_M=64}, \texttt{BLOCK\_N=64}, \texttt{BLOCK\_K=32}, \texttt{BLOCK\_R=32}) and warp counts (\texttt{num\_warps=4}) for arbitrary rank configurations, whereas static CUDA kernels without runtime tile autotuning can suffer catastrophic execution penalties. The \neuralops{} dispatch engine actively monitors execution profiles, routing requests to the autotuned Triton backend ($0.041\,\text{ms}$) to avoid the un-tuned CUDA fallback penalty.

% ==============================================================================
% D.6 Engineering Guidelines for Cross-Silicon Deployment
% ==============================================================================
\subsection{Engineering Guidelines for Cross-Silicon Deployment}
\label{sec:appendix_engineering_guidelines}

Based on the empirical telemetry and microarchitectural dissections presented in this Appendix, we formalize five structural rules for cross-silicon LLM kernel engineering:
\begin{enumerate}
    \item \textbf{Enforce Launch Bounds on Ada Lovelace}: Always decorate custom CUDA kernels with \texttt{\_\_launch\_bounds\_\_(128, 2)} or \texttt{\_\_launch\_bounds\_\_(256, 2)}. Restricting register allocation to $\le 64$ registers per thread avoids the $50\%$ occupancy cliff and guarantees adequate warp scheduling headroom to hide GDDR6X latency.
    \item \textbf{Pad Shared Memory Strides to Prime Offsets}: Avoid power-of-two leading matrix dimensions in shared memory tiles. Incorporate an 8-word padding stride ($\text{Stride} = \text{Dim} + 8$) to eliminate 32-way shared memory bank collisions during column sweeps.
    \item \textbf{Enforce 32-Byte and 512-Byte Alignment on Ascend}: Ensure that all tensor pointers, chunk lengths, and head dimensions dispatched to Huawei Ascend NPUs are strict multiples of 32 bytes (minimum hardware alignment) and ideally 512 bytes (DMA burst efficiency limit).
    \item \textbf{Preserve 16-Element Multiples for Ascend Cube Engagement}: Never deploy non-power-of-two head projections ($d_k \ne 16k$) to Ascend NPUs without upfront padding. Forcing matrix multiplication onto the Vector Unit incurs an immediate $10\times$ throughput penalty.
    \item \textbf{Implement Identity Copies on Ascend Triton Accumulators to Prevent DMA Race Conditions}: In chunked linear attention and recurrent state update loops ($\mat{S}_c = \mat{S}_{c-1} + \mat{K}_c^\top \mat{V}_c$) running under the \texttt{triton\_ascend} backend, developers must enforce the identity assignment \texttt{k\_compute = k\_buf0 + 0.0} prior to invoking \texttt{tl.dot}. Microarchitecturally, the Cube Unit loads matrix tiles asynchronously via MTE1 from the Unified Buffer into internal L0A/L0B systolic registers. In a three-stage double-buffered pipeline, while tile $c$ is being staged, MTE2 initiates asynchronous burst DMA from HBM to stage tile $c+1$ into the alternate UB partition. Because the BiSheng compiler's liveness analysis operates on register SSA lifespans rather than physical DMA transfer barriers, it prematurely marks the memory slice of \texttt{k\_buf0} as dead once the transposed reference is consumed. Consequently, inbound MTE2 transfers overwrite the trailing memory blocks of $\mat{K}_c$ in UB before systolic staging into L0A is complete, inducing silent $\text{NaN}/\pm\infty$ numerical divergence. The addition of floating-point zero creates a distinct Vector Unit instruction (\texttt{VADD}) and an independent SSA live range, forcing the compiler allocator to maintain physical memory exclusivity until the systolic operation commits and synchronizes across hardware pipeline barriers.
\end{enumerate}
